\section{Experiments, replay, and what the evidence does not say}
\label{sec:experiments}

\subsection{Execution protocol}
The recorded run used Python 3.13.5 on Linux, with seed \code{20260915}. All
commands ran with \code{python -S}, which disables automatic site-package
loading. The implementation imports only standard-library modules. The source
and complete generated certificates are included in the archive, together with
the run's summary and console logs.

The experiments distinguish a \emph{model case}, a \emph{start-state check},
and a \emph{stored certificate record}. A game may have both a nonempty winning
region and a nonempty losing region, producing two stored records for one
arena. Every vertex is checked as an initial vertex during the experiment, but
replay stores one representative start per region. Thus the record count is
not a count of distinct theorem problems and certainly not a count of
Lean-kernel proofs.

\begin{table}[ht]
\centering\small
\begin{tabular}{@{}lrrr@{}}
\toprule
Family & Model cases & Positive records & Negative records \\
\midrule
Pushdown: non-growing & 500 & 199 & 301 \\
Pushdown: general normalized & 250 & 70 & 180 \\
Pushdown: regular targets & 100 & 56 & 44 \\
Pushdown: compressed runs & 31 & 31 & 0 \\
\midrule
Games: exhaustive sizes 1--3 & 22,100 & 15,388 & 11,260 \\
Games: random sizes 4--7 & 200 & 153 & 114 \\
\midrule
Chains: acyclic & 200 & 200 & 0 \\
Chains: cyclic/Cramer oracle & 150 & 150 & 0 \\
Chains: geometric & 60 & 60 & 0 \\
Chains: reachable traps & 60 & 0 & 60 \\
Chains: arbitrary four-state & 300 & 250 & 50 \\
\bottomrule
\end{tabular}
\caption{Recorded cases and certificate outcomes. A negative record is an
expected mathematical result, not a failed solver invocation. Game records
count nonempty regions, not arenas.}
\label{tab:counts}
\end{table}

Across these datasets the archive contains 28,566 certificate records:
881 pushdown records, 26,915 game-region records, and 770 chain records. All
replay successfully. This total describes the archive's transport workload;
it is not a pooled success-rate statistic. No results from the repository's
eighteen historical proposals are included in these numbers.

\subsection{Pushdown reference computations and their limits}
The 500 non-growing systems use rules of replacement length at most one. Their
reachable configuration spaces are finite, so concrete breadth-first search
from the initial stack is a complete oracle. Its verdict agrees with the
summary solver in every case; positive certificates are also expanded into
actual runs and checked against the original rules.

The 250 general systems allow pushes of length two. A separate breadth-first
search with maximum stack height six finds 70 positive instances and completely
exhausts 141 negative instances. In 39 further cases it encounters transitions
that exceed its height limit and must report \emph{unknown}. The summary worker
returns exclusion certificates for those cases. They pass the independent
closure audit, but the bounded concrete oracle does not independently settle
them. The run records this limitation rather than counting truncated searches
as corroborating negative answers.

The regular-target suite generates non-growing original systems and random
three-state target DFAs. It compares a direct predicate-on-stack breadth-first
search with the compiled empty-stack query. All 100 verdicts agree. The stored
certificate records concern the compiled systems; reproduction of the
experiment reconstructs the original target fixtures and performs the
source/compiled comparison. Certificate replay alone does not re-prove the
compiler bridge.

The compressed family covers depths zero through thirty. All represented
lengths equal $2^{d+1}-1$ and all positive DAGs have $d+1$ nodes. Cases through
depth ten are expanded and replayed as concrete traces. Larger cases validate
the shared certificate and its length recurrence without materializing the
trace. The maximum example's length is therefore \emph{represented}, not an
observed count of executed program steps.

\subsection{An exhaustive finite-game corpus}
For each $n\in\{1,2,3\}$, the experiment enumerates every nonempty successor
subset at each labelled vertex, every binary ownership map, and every accepting
subset. The number of arenas is
\[
 \sum_{n=1}^{3}(2^n-1)^n\,2^n\,2^n
   =4+144+21,952=22,100.
\]
Self-loops are allowed. Isomorphic labelled arenas are not deduplicated; the
claim is exhaustive coverage of this explicit finite input domain, not a count
of nonisomorphic games. An additional 200 arenas have four to seven vertices
and at most three distinct outgoing edges per vertex.

The reference oracle does not use attractors or the solver's rank construction.
For each player's positional policy, it fixes that player's choices and leaves
all opposing choices available. It computes transitive closure by a separate
Floyd-style routine. For a protagonist policy it detects reachable cycles in
the nonaccepting subgraph. For an antagonist policy it detects reachable cycles
containing an accepting vertex. It unions the respective winning sets over
all positional policies and checks that they partition the arena. The search's
winning and losing regions agree with this oracle on every arena.

The experiment additionally checks every initial vertex with its region
certificate: 67,260 start-state replays, comprising 40,250 winning and 27,010
losing starts. The persistent artifact stores only one representative start
per nonempty region. This difference explains why game checks, records, and
arenas have different counts.

This is substantially stronger evidence than a handful of successful examples,
but it is still finite testing. It shares the standard positional-strategy
semantics with the solver and does not replace the soundness and completeness
arguments or their future Lean formalization.

\subsection{Exact-chain oracles}
The 200 acyclic chains are evaluated by backward finite-path dynamic
programming. The 150 small cyclic chains have values checked by Cramer's rule,
with determinants computed by Laplace expansion rather than elimination. The
60 geometric chains are checked against the closed forms $\E\tau=d$,
expected cost $3d/2$, and terminal value one for exit probability $1/d$.
Thus 410 positive cases have numerical values compared against a genuinely
different computation or analytic formula.

The 60 explicit trap examples have a known entry probability $1/d$, and the
returned lower bound agrees exactly. The 300 arbitrary four-state chains use
an oracle that enumerates nonempty closed subsets of reachable nonterminal
states rather than computing SCCs. It verifies the almost-sure classification,
not every numerical field of the positive certificate. All 250 positive and
50 negative classifications agree. Their numerical residuals and path/closure
conditions are still checked by the replay module.

This separation is intentional. ``Independent oracle'' should identify the
claim independently checked: exact values in the first three groups,
closed-class probability in the fourth, and almost-sure classification in the
fifth. It should not be stretched into a statement that every output field was
computed twice by unrelated implementations.

\subsection{Negative controls and transport tests}
The 36 named unit tests all pass. They include successful examples as well as
invalid-input and certificate mutations, so 36 is not a count of rejected
corruptions. Representative controls reject a self-referential pop proof,
missing pop bases or binary closure consequences, the wrong pushdown target,
an old policy under changed ownership, an omitted opponent self-loop, and a
counterrank that cannot pay for an accepting visit.

Probability controls reject nonunit rows, floats, changed transition
probabilities, omitted positive-mass successors, a path through a zero-mass
edge, a purported trap containing a terminal, and an arbitrary harmonic value
with no properness potential. The tests also confirm that an unreachable
nonterminal trap does not invalidate absorption from the actual start.
Transport controls reject duplicate JSON keys, zero rational denominators, and
Python booleans used as integer state identifiers.

These tests exercise real boundaries in the implemented formats. They are not
a security audit. The research decoder has an input-size cap, duplicate-key
rejection, and rational integer-size checks, but it has not been hardened
against every hostile nesting pattern, pathological dimension, or resource
exhaustion attack. Unexpected malformed shapes may raise ordinary Python
exceptions. A service implementation needs a separately bounded, total decoder
and process isolation; none of those are claimed as delivered here.

\subsection{Search-free replay and performance interpretation}
The replay command loads the separately encoded problem, decodes its
certificate, calls the appropriate checker, and compares the result with the
recorded expected outcome. The checker's own conclusion is recomputed; the
\code{expected} field is only a regression assertion. Before and after the
run, the script checks that the three search modules have not been imported.
It prints:
\begin{lstlisting}
Replayed 28566 records; no search modules imported.
reachable: 356             unreachable: 525
winning: 15541             losing: 11374
almost_sure: 660           not_almost_sure: 110
\end{lstlisting}

Search and replay still share the input model definitions, the codec, and
Python's rational arithmetic implementation. This is more separation than
calling the search again during replay, but less than an independently
verified checker. Likewise, excluding site packages proves that these runs did
not require an installed numerical solver; it does not prove that the code is
correct or malicious-input safe.

The summary records internal search, check, and game-oracle timings. The
complete experiment loop took approximately 1.62 seconds in the recorded run,
excluding JSON serialization, process startup, unit tests, and the later
standalone replay. This is a single small-model run, with different numbers of
search and replay calls, not a statistical performance study. The timings are
retained for reproducibility and troubleshooting, not as a speed claim.

No Lean executable was available in this environment. No Lean source from this
package was elaborated, no model theorem was checked by the Lean kernel, and no
comparison with \code{grind}, Aesop, LeanHammer, or another tactic was run.
The historical core Lean elaborations in Forge's README belong to that
repository and are not credited to this work. An end-to-end capability gain in
Lean remains a future measurement, although the Python workers already decide
their explicitly specified model-level fragments.
